\section{The infinite coset complex}
\label{sec:geometry}

Put $L=\CC(U;(H_i)_{i\in S})$. We will prove
\begin{equation}\label{eq:L-acyclic}
 \wtH_j(L;\F)=0\qquad(0\le j<n-1).
\end{equation}
This is the acyclicity needed to compare group homology with the homology
of the finite complexes. We first identify the finite stabilizers, then
realize $L$ inside an affine building and compute the relative chains of
a filtration. Connectedness will also show that the $x_i$ generate $U$.

\subsection{The finite local groups}

Regard $S$ as the edges of the directed cycle
$1\to2\to\cdots\to N\to1$, with edge $0$ equal to $N\to1$.
For $J\subsetneq S$, put $U_J=\langle x_j:j\in J\rangle$ and
$U_{\varnothing}=\{1\}$. Between distinct indices $a,b$ there is a unique
simple directed path in the cycle. Set
\[
 V_{ab}(J)=
 \begin{cases}
 k,&\text{if that path lies in $J$ and avoids edge $0$},\\
 kt,&\text{if that path lies in $J$ and uses edge $0$},\\
 0,&\text{otherwise}.
 \end{cases}
\]

\Needspace{12\baselineskip}
\begin{lemma}\label{lem:patterns}
For every proper $J\subsetneq S$,
\begin{equation}\label{eq:pattern}
 U_J=\{g:g_{aa}=1,\ g_{ab}\in V_{ab}(J)\text{ for }a\ne b\}.
\end{equation}
These are finite groups, and
\begin{equation}\label{eq:intersections}
 U_J\cap U_K=U_{J\cap K},\qquad
 \bigcap_{i\in A}H_i=U_{S\setminus A}\quad(\varnothing\ne A\subseteq S).
\end{equation}
For $m\ge2$, reduction modulo $t^m$ is injective on all $U_J$ and preserves
all these intersections. In particular,
\begin{equation}\label{eq:M}
 |H_i|=M:=q^{\binom N2}\quad(i\in S),\qquad |U_{\{i\}}|=q.
\end{equation}
\end{lemma}

\begin{proof}
Deleting an edge outside $J$ cuts the cycle into a linear order. In this
order the right side of \eqref{eq:pattern} is a product of unitriangular
groups, one for each path component of $J$. Factors of $t$ correspond to
conjugation by a diagonal matrix over $k[t,t^{-1}]$. In particular the
displayed matrices form a group.

Because $q$ is prime, powers of $x_j$ provide every coefficient in its
simple root group. Iterated commutators
\[
 [e_{ab}(f),e_{bc}(g)]=e_{ac}(fg)
 \qquad(a,b,c\text{ distinct})
\]
provide every allowed path entry. Unitriangular elimination in the cut
order then gives every matrix on the right of \eqref{eq:pattern}.
This proves the equality. The path descriptions also show the
intersection identities. Finally, every entry has degree at most one;
reduction modulo $t^m$, for $m\ge2$, neither identifies distinct such
entries nor changes their allowed subspaces. The same intersection
identities therefore hold after reduction.
\end{proof}

\subsection{The twin-building model}

The matrix calculation identifies the stabilizers we need. To obtain
acyclicity, we now place the coset complex in the classical affine
twin building of $\SL_N(k[t,t^{-1}])$. We use the lattice construction
of Abramenko--Van Maldeghem
\cite[Section~4, Proposition~4.4 and Remark~4.6]{AVM2001};
see also \cite[Lemma~5]{Abramenko} and
\cite{AbramenkoBrown}. Its two halves are the
affine buildings at $0$ and at infinity, both of type $\widetilde A_n$
and dimension $n$. Their chambers have one vertex of each type in $S$.
The stabilizers of opposite base chambers are
\[
 B_+=\{g\in\SL_N(k[t]):g(0)\text{ is upper triangular}\},
\]
and the analogous group $B_-$ over $k[t^{-1}]$ with lower triangular
reduction at infinity. They satisfy $B_+=U\rtimes T$ and
$B_+\cap B_-=T$, where $T$ is the constant diagonal torus.

Fix the positive base chamber $c_+$. A chamber $c$ in the negative
building has a \emph{codistance} $\delta^*(c)=\delta^*(c_+,c)\in W$,
where $(W,S)$ is the affine Coxeter system. For two chambers in the same
half, their Weyl distance is the element of $W$ represented by the types
of a minimal gallery between them. Write $\ell$ for word length with respect to $S$, and
$W_J=\langle J\rangle$ for $J\subseteq S$.
The chambers of codistance $1$ are exactly the
$B_+$-orbit of the negative base chamber. We use the convention that an
$s$-panel varies the vertex of type $s$. If $\delta^*(c)=w$, the
codistances in this panel belong to $\{w,ws\}$; the value $ws$ is attained,
and, when $\ell(ws)<\ell(w)$, every chamber in the panel other than $c$
has codistance $ws$. These are the twin-building axioms with the positive
argument fixed \cite[Definition~1.1]{AVM2001}.
Every proper $J\subsetneq S$ is spherical:
a $J$-residue, obtained by varying only types in $J$, is a finite spherical
building whose irreducible factors have type $A$ over $k$. In the lattice
model, fixing the other types fixes a partial lattice chain; completing
that chain amounts to choosing flags in its successive finite-dimensional
$k$-vector-space quotients. This gives the asserted type-$A$ factors.

Let $L'$ be the subcomplex consisting of all codistance-$1$ chambers and
their faces. The preceding description shows that $U$ acts simply
transitively on its chambers. We check their vertex stabilizers.
In the standard lattice model at infinity, vertices are homothety classes
of $k[[t^{-1}]]$-lattices. Choose the negative base chamber with vertices
represented by exponent vectors
\[
 (0,\ldots,0)\quad\text{(type $0$)},\qquad
 (\underbrace{0,\ldots,0}_{i},\underbrace{1,\ldots,1}_{N-i})
 \quad\text{(type $i$)}.
\]
For a vertex represented by exponents $(b_a)$, a polynomial matrix in
its stabilizer satisfies
\begin{equation}\label{eq:lattice-bounds}
 \deg g_{ab}\le b_a-b_b.
\end{equation}
This is the condition for preserving the lattice generated by $t^{b_a}e_a$
over $k[[t^{-1}]]$; because the determinant is $1$, preservation of its
homothety class implies preservation of the lattice itself. Intersecting
these bounds with $U$ gives upper unitriangular constant matrices at the
type-$0$ vertex. At type
$i\ne0$ they give two upper unitriangular constant diagonal blocks, cut
after coordinate $i$, a zero upper-right block, and arbitrary entries
in $kt$ in the lower-left block. By \cref{lem:patterns}, these stabilizers
are exactly $H_0,H_i$, respectively. Sending $uH_i$ to the corresponding
vertex identifies
\begin{equation}\label{eq:L-opposite}
 L=\CC(U;(H_i))\cong L'.
\end{equation}
This uses the common-representative definition of faces: a face of $L'$
lies in one of its chambers, and every such chamber is a unique translate
of the base chamber.

\subsection{A filtration by codistance}

We have identified $L$ with an opposition subcomplex. The next step is
to compare it with the whole negative building by adding chambers in
codistance order. The homology of each attachment is controlled by
spherical opposition complexes, for which we use the following input.
In a
spherical building of type $A_r$ over $k$, let $O$ be the subcomplex formed
by chambers opposite a fixed chamber and their faces. For
$q\ge2^{r-1}$,
\begin{equation}\label{eq:spherical-opposition}
 \wtH_j(O;\F)=0\qquad(j<r-1).
\end{equation}
This follows from the integral cone function in
\cite[Corollary~6.10]{OVB}, whose chain contraction identity remains valid
after reduction modulo two. For $r=1$, the required assertion is just
nonemptiness of the opposition complex. For a product of spherical buildings, the
opposition complex is the simplicial join of the opposition complexes
of the factors, and the augmented chain formula for joins gives the
same vanishing below its dimension. We henceforth require $q\ge2^{n-1}$,
which covers every proper residue of our affine building.

\begin{proposition}\label{prop:L-acyclic}
The complex $L$ satisfies \eqref{eq:L-acyclic}.
\end{proposition}

\begin{proof}
Starting with $L$, add all chambers of the negative building in increasing
order of the length of their codistance. Order the finitely many
codistance values of each length arbitrarily. Fix a value $w\ne1$, and put
\[
 D=\{s\in S:\ell(ws)<\ell(w)\},\qquad J=S\setminus D.
\]
The descent set $D$ is nonempty, so $J$ is proper and $W_J$ is finite.
The element $w$ is the minimal representative of $wW_J$; in particular
$\ell(wv)=\ell(w)+\ell(v)$ for $v\in W_J$.

Group the chambers of codistance $w$ by their $J$-residues. Such a residue
is specified by the common face $F$ whose vertex types are $D$. Its
varying part is a spherical building of type $J$.
We~claim that the codistance-$w$ chambers in this residue are precisely
the chambers opposite a single chamber $c_{\max}$ there. Indeed, follow
a reduced expression for the longest element $w_J$ from any chamber of
codistance $w$, choosing at each step the codistance increase permitted
by the panel axiom. This reaches $c_{\max}$ of codistance $ww_J$.
From $c_{\max}$, any minimal gallery within the residue has types forming
a reduced word in $W_J$. Every step decreases codistance length, since
\[
 \ell(ww_Jv)=\ell(w)+\ell(w_J)-\ell(v)\qquad(v\in W_J).
\]
The descent panel axiom thus forces the endpoint at Weyl distance $v$
from $c_{\max}$ to have codistance $ww_Jv$. It has codistance $w$ exactly
when $v=w_J$, which proves the claim. It also proves that no chamber
in this residue has smaller codistance length or a different codistance
of the same length. The same formula proves uniqueness of $c_{\max}$:
a chamber of codistance $ww_J$ has Weyl distance $1$ from it.

Let $O$ denote this opposition complex. The chambers in the residue
now being added, together with their faces, form $F*O$. Its intersection
with the previously added subcomplex is exactly $\partial F*O$.
For one inclusion, a face missing a vertex of type $s\in D$ lies in the
corresponding panel of a codistance-$w$ chamber, and hence in a chamber
of codistance $ws$, which was added earlier. Conversely, a face
containing all of $F$ can lie only in chambers of the same $J$-residue,
so the preceding paragraph shows that it was not present earlier.
Distinct blocks for this $w$ meet only in their attaching loci, since
a face outside such a locus determines its face of types $D$.

The relative chain complex of a block is particularly simple:
\begin{equation}\label{eq:relative-chains}
 C_k(F*O,\partial F*O;\F)
 \cong \wtC_{k-|D|}(O;\F),
 \qquad
 H_k(F*O,\partial F*O;\F)
 \cong \wtH_{k-|D|}(O;\F).
\end{equation}
Only faces containing all vertices of $F$ survive in relative chains;
deleting one of those vertices gives zero, leaving the augmented
boundary in $O$, as illustrated in Figure~\ref{fig:codistance-attachment}.
\begin{figure}[tbp]
\centering
\begin{tikzpicture}[
  every node/.style={font=\small,align=center},
  box/.style={draw,rounded corners=2pt,inner sep=6pt,text width=5.0cm},
  old/.style={box,draw=black!55,fill=black!5},
  new/.style={box,draw=blue!50!black,fill=blue!4},
  arrow/.style={-{Latex[length=2mm]},semithick}
]
\node[box,text width=6.3cm] (block) at (0,2.8)
  {A face $\tau$ of the block $F*O$\\
   $F$ has $|D|$ vertices; $\dim O=|J|-1$};
\node[old] (old) at (-3.3,1.0)
  {$\tau$ omits a vertex of $F$\\
   $\tau\in\partial F*O$: already present};
\node[new] (survives) at (3.3,1.0)
  {$\tau=F\cup\sigma$, with $\sigma\in O$\\
   Survives in relative chains; $\sigma=\varnothing$ is allowed};
\draw[arrow] (block.south west) -- (old.north);
\draw[arrow] (block.south east) -- (survives.north);
\node[old] (deletef) at (-3.3,-1.1)
  {A vertex of $F$ is deleted\\
   The resulting face is zero in relative chains};
\node[new] (deleteo) at (3.3,-1.1)
  {A vertex $v$ of $\sigma$ is deleted\\
   The resulting face is $F\cup(\sigma\setminus\{v\})$};
\draw[arrow] (survives.south west) -- (deletef.north east);
\draw[arrow] (survives.south) -- (deleteo.north);
\node at (0,-2.35)
  {$C_k(F*O,\partial F*O;\F)
    \cong \wtC_{k-|D|}(O;\F)$};
\end{tikzpicture}
\caption{One block of the codistance filtration, shown as a schematic
of its faces and boundary operations. The new relative faces contain
every vertex of $F$. Deleting an $F$-vertex lands in the old subcomplex;
the remaining boundary is the augmented boundary of $O$, shifted by
$|D|$. The complex $O$ is not assumed to be contractible.}
\label{fig:codistance-attachment}
\end{figure}

Since $\dim O=|J|-1$, \eqref{eq:spherical-opposition}
shows that these relative groups vanish for
\[
 k<|D|+|J|-1=|S|-1=n.
\]
If $J=\varnothing$, interpret $O$ as the complex consisting of the empty
face; the pair is $(F,\partial F)$ and the same conclusion holds.

Relative chains for the whole step are the direct sum of the block
relative chains. The long exact sequence of the pair therefore shows
that every step induces an isomorphism in reduced homology in degrees
$j<n-1$. The filtration exhausts the negative affine building, which is
contractible \cite[Theorem~II.10A.4 and Remark~II.10A.6]{BH1999},
and simplicial homology commutes with this increasing union. Thus
$\wtH_j(L;\F)=0$ in precisely the asserted range.
\end{proof}

\begin{corollary}\label{cor:generation}
The simple generators $x_i$, $i\in S$, generate $U$.
\end{corollary}

\begin{proof}
By \cref{prop:L-acyclic}, $L$ is connected. In a connected pure complex,
any two chambers can be joined by a sequence of chambers in which
successive chambers share a vertex. If the chambers labeled $u,v\in U$
share a vertex of type $i$, then $uH_i=vH_i$, so $u^{-1}v\in H_i$.
The chamber action is simply transitive, so following such a sequence
from the base chamber shows that $U=\langle H_i:i\in S\rangle$.
Each $H_i$ is generated by a subset of the $x_j$.
\end{proof}
